%%%%%%%%%%%%%%%%%%%%%%%%%%%%%%%%%%%%%%%%%%%%%%%%%%%%%
%%   Atencao! A lista nao ordena automaticamente   %%
%%%%%%%%%%%%%%%%%%%%%%%%%%%%%%%%%%%%%%%%%%%%%%%%%%%%%

\begin{simbolos} \itemsep -1pt
	\item[$ A $] Área de uma bacia hidrográfica (km$^2$)
	\item[$ c $] Pré-fator de uma lei de escala $y = c\,x^{\beta}$
	\item[$ D^{*} $] Limiar de densidade populacional do CCA (hab/km$^2$)
	\item[$ N $] Número de maiores manchas usadas no estimador de Hill
	\item[$ n $] Tamanho de uma amostra
	\item[$ P $] Fração da área de terra coberta pelas bacias principais
	\item[$ p $] Valor-$p$ de um teste de hipótese
	\item[$ p(x) $] Densidade de probabilidade do tamanho $x$
	\item[$ R $] Raio de aceitação entre a foz e o ponto de coleta (km)
	\item[$ R_{\mathrm{maioria}} $] Menor raio que contém ao menos metade dos pontos de coleta
	\item[$ r $] Posto (\textit{rank}) de uma mancha urbana; coeficiente de correlação de Pearson
	\item[$ S(r) $] População da mancha urbana de posto $r$
	\item[$ x_{\min} $] Limite inferior da cauda ajustada por lei de potência
	\item[$ \alpha $] Expoente de Zipf na forma \textit{rank-size}, $S(r) \propto r^{-\alpha}$
	\item[$ \alpha_N $] Expoente de Zipf estimado por Hill sobre as $N$ maiores manchas
	\item[$ \beta $] Expoente de uma lei de escala (alometria)
	\item[$ \gamma $] Expoente da densidade de uma lei de potência, $p(x) \propto x^{-\gamma}$
	\item[$ \Delta $] Desvio máximo $\max_N |\alpha_N - 1|$ usado na calibração
	\item[$ \ell $] Distância máxima de ligação entre pixels do CCA (m)
	\item[$ \ell_{\min} $] Menor $\ell$ admitido na calibração (m)
	\item[$ \rho_s $] Coeficiente de correlação de postos de Spearman
	\item[$ \zeta $] Expoente de cauda estimado por Hill, $\zeta = 1/\alpha$
\end{simbolos}
